\section{Zadanie 8}
\autor{Dominik Korsa}
Rzucamy monetą, gdzie orzeł wypada z prawdopodobieństwem \(p\). Każdy rzut trwa
czas będący zmienną losową z rozkładu wykładniczego z parametrem 1.
Pokaż, że łączny czas oczekiwania na orła jest zmienną losową z rozkładu
wykładniczego o parametrze \(p\).

\subsection*{Rozwiązanie}
Niech \(X\) to zmienna losowa określająca liczbę rzutów do wypadnięcia orła.
\(X\) ma rozkład geometryczny z parametrem \(p\).

Niech \(Y_i\) to czas trwania \(i\)-tego rzutu.
Wtedy \(Y = Y_1 + \dots + Y_X\) to czas oczekiwania na orła.

Rozważmy gęstość w zależności od wartości \(X\):
\begin{align*}
	f_Y(y)
	 & = \sum_{x = 1}^{\infty} f_{Y \mid X = x}(y) \cdot P(X = x) \\
\end{align*}

Pod warunkiem \(X=x\) jest tak, że \(Y = Y_1 + \cdots + Y_x\).
\(Y\) jest sumą \(x\) niezależnych zmiennych o rozkładzie wykładniczym z parametrem
\(1\), więc z \hyperref[set09.task05]{zadania 5} \(Y\) ma rozkład
\(\Gamma(x, 1)\). Ze wstępu do zestawu znamy gęstość zmiennej o takim rozkładzie:
\[
	f_{Y \mid X = x}(y) = \frac{1}{\Gamma(x)} y^{x-1} e^{-y}
\]
A dodatkowo wiemy, że \(X\) ma rozkład geometryczny, więc
\(P(X = x) = (1-p)^{x-1} p\).
Mamy:
\begin{align*}
	f_Y(y)
	 & = \sum_{x = 1}^{\infty} \frac{1}{\Gamma(x)} y^{x-1} e^{-y} \cdot (1-p)^{x-1} p \\
\end{align*}
Wyciągnijmy stałe czynniki przed sumę i uprośćmy.
Wykorzystajmy dodatkowo fakt, że skoro \(x \in \mathbb{N}_+\),
to \(\Gamma(x) = (x-1)!\).
W drugiej linijce przesuwamy iterator w sumie:
\begin{align*}
	f_Y(y)
	 & = p e^{-y} \sum_{x = 1}^{\infty} \frac{1}{(x-1)!} y^{x-1} \cdot (1-p)^{x-1} \\
	 & = p e^{-y} \sum_{x = 0}^{\infty} \frac{(y-yp)^x}{x!}
\end{align*}
W tym momencie udajemy, że mieliśmy porządną analizę, rozpoznajemy tutaj wzór
\[e^z = \sum_{k=0}^{\infty} \frac{z^k}{k!}\]
i otrzymujemy:
\[
	f_Y(y) = p e^{-y} e^{y-yp} = p e^{-py}.
\]
To jest dokładnie gęstość zmiennej losowej z rozkładu wykładniczego z parametrem
\(p\), więc \(Y\) jest taką zmienną losową, co kończy dowód.
